\documentclass[11pt]{article}
\usepackage[margin=1in]{geometry}
\usepackage{amsmath, amssymb}

\title{Exercise 3.18: $v_\pi(s)$ from the Backup Diagram}
\author{Sutton \& Barto, Section 3.5}
\date{}

\begin{document}
\maketitle

\section*{Problem}

The value of a state depends on the values of the actions possible in that
state and on how likely each action is to be taken under the current policy.
The backup diagram is rooted at state $s$ and branches to possible actions.

Give the equation for $v_\pi(s)$ in terms of $q_\pi(s, a)$, given $S_t = s$.
First with an expectation conditioned on following $\pi$, then written out
explicitly with no expected-value notation.

\section*{Solution}

\subsection*{With expectation notation}

\[
    v_\pi(s) = \mathbb{E}_\pi\!\left[\, q_\pi(s, A_t) \mid S_t = s \,\right]
\]

The expectation is over the action $A_t$ selected by $\pi$ in state $s$.

\subsection*{Written out explicitly}

\[
    \boxed{v_\pi(s) = \sum_{a} \pi(a \mid s)\; q_\pi(s, a)}
\]

This corresponds to the upper half of the $v_\pi$ backup diagram (Figure~3.4,
left): from the state node, we branch to each action $a$ with probability
$\pi(a \mid s)$ and read off $q_\pi(s, a)$ at the leaves.

\end{document}
